\section*{Resumen}

Este trabajo presenta un sistema de Inteligencia Artificial capaz de tomar decisiones de descarte en Riichi Mahjong, un juego tradicional japonés de información imperfecta caracterizado por un elevado nivel de complejidad estratégica.

Para ello, se desarrolla un modelo basado en la arquitectura Transformer entrenado mediante aprendizaje por imitación a partir de un conjunto de datos compuesto por partidas reales obtenidas de la plataforma Tenhou.net. El rendimiento del modelo se evalúa utilizando métricas convencionales de precisión, analizando además las implicaciones de los resultados obtenidos.

Con el objetivo de complementar estas métricas tradicionales, se propone en este trabajo una metodología de evaluación denominada \textit{Precisión Estratégica}, basada en categorías definidas mediante Shanten y Ukeire. Este enfoque permite analizar la eficiencia de las decisiones generadas por el modelo más allá de la coincidencia exacta con la acción observada.

Los resultados experimentales muestran que el modelo entrenado con 500.000 estados alcanza una precisión exacta del 65,64\%, una precisión top-3 del 92,47\% y una precisión estratégica del 94,01\% sobre un conjunto reservado de 52.563 estados. En comparación con el checkpoint identificado como 100k, el segundo modelo mejora todas las métricas predictivas y reduce las decisiones que empeoran el Shanten del 5,03\% al 1,59\%. Estos resultados indican que el modelo aprende regularidades relevantes para la selección de descartes, aunque la evaluación estratégica queda limitada a los aspectos representados mediante Shanten y Ukeire.

\vspace{0.5cm}

\textbf{Palabras clave:}

Riichi Mahjong,
Inteligencia Artificial,
Transformers,
Aprendizaje por Imitación,
Precisión Estratégica,
Evaluación de Modelos,
Juegos de Información Imperfecta.
